
\chapter{Introduction}
\label{chap:introduction}
\thispagestyle{empty}  % first page of Introduction: no page number; next page = Arabic 2

Images acquired in high-altitude or extreme-weather regions such as fog, smog, dust storms, haze or snow are often characterised by low contrast, reduced visibility, colour distortion and loss of fine structural details. These degradations create a distribution shift between clean datasets and real-world imagery, leading to poor performance of transformer-based vision systems like Vision Transformers (ViT)~\cite{dosovitskiy2020vit} and multimodal encoders such as CLIP~\cite{radford2021clip} in adverse environmental conditions. Most modern vision backbones are tailored for clear, well-lit scenes; they endure significant performance degradation when deployed on aerial platforms, agricultural monitoring systems, autonomous vehicles, or environmental sensing applications under challenging weather conditions.

\noindent The proliferation of drone-based monitoring systems in agriculture, disaster management, environmental surveillance and autonomous navigation has emphasised the requirement for robust visual perception in diverse environmental conditions. Atmospheric phenomena such as fog, haze, smoke, snow, sandstorms, and dust accumulation degrade image formation processes by limiting visibility and contrast through scattering effects, which often leads to drastic degradation of recognition and classification performance in downstream tasks. Large-scale datasets covering such environmental variations are rare, difficult to collect and expensive to annotate, despite their practical importance. This limitation has led to growing interest in synthetic data generation as a scalable alternative for studying and mitigating environmental distribution shifts.

\noindent Recent progress in diffusion-based generative models~\cite{ho2020ddpm} has allowed for the generation of highly realistic images by iteratively denoising the image conditioned on textual, visual, and structural information. Compared to traditional augmentation methods, diffusion models offer fine-grained control of scene content and environmental appearance, making them ideal candidates for realistic atmospheric simulation. Through mechanisms such as cross-attention, adapter-based conditioning, latent-space guidance, and image-conditioned generation, diffusion models can produce visually coherent images while maintaining semantic structure, thereby offering an effective framework for synthetic dataset generation.

\noindent A large body of work has investigated style-transfer and image-conditioned diffusion frameworks for controllable image synthesis. Methods such as StyleShot~\cite{styleshot2024}, CSGO~\cite{csgo2024}, InstantStyle~\cite{wang2024instantstyle}, and other adapter-based approaches have shown success in transferring colour distributions, illumination properties, textures, and environmental appearance from reference images to generated scenes. Although these techniques often produce visually realistic atmospheric effects, the transferred style is frequently biased towards background regions and global image appearance, leaving foreground objects largely unaffected. Atmospheric effects such as haze, fog, snow, or dust are typically represented as colour changes, texture overlays, and scene mood shifts, while object boundaries, semantic structure, and class-discriminative visual cues are mostly maintained. Consequently, the resulting images can look very different to human observers while still preserving the semantic information used by modern recognition systems.

\noindent This discrepancy underscores an important distinction between perceptual realism and representation-level distribution shift. A visually realistic atmospheric transformation does not always imply a meaningful change in the feature representations learned by modern vision models. In many cases, the appearance of the environment can vary greatly, but foreground objects and the corresponding high-confidence predictions of large vision-language models  can still be recognised reasonably well. This poses a critical challenge for atmospheric style transfer and raises important questions about the effectiveness of appearance-based approaches in generating truly shifted domains.

\noindent In addition to the challenge of generating synthetic images, atmospheric degradations pose further challenges to deployed recognition systems. Severe weather conditions suppress high-frequency information, blur object boundaries, alter feature distributions, and reduce semantic clarity. These effects motivate the development of adaptation strategies that can improve model robustness when source-domain data is unavailable or inaccessible. Recent advances in vision-language models, synthetic source reconstruction, bridge-domain alignment, and frequency-aware adaptation mechanisms~\cite{guibas2021afno} have pointed to promising directions for addressing such challenges. Understanding the link between atmospheric domain generation and model adaptation is therefore essential for building perception systems that can reliably operate in adverse environmental conditions.

\noindent The controllability of diffusion models has also fostered the development of guidance mechanisms for directly controlling the denoising trajectory. Beyond conventional text conditioning, perceptual style descriptors, latent-space correction strategies, contrastive objectives, and optimisation-inspired guidance techniques provide additional control over generated outputs. These approaches offer the potential for more robust atmospheric transformations while preserving scene consistency, enabling a more systematic investigation of environmental distribution shifts and their impact on modern vision systems.

\noindent While appearance-based stylisation methods are effective in modifying visual attributes related to adverse weather conditions, many existing methods do not explicitly model the physical processes underlying atmospheric image formation. Real-world phenomena such as haze, fog, smoke, and dust are caused by depth-dependent scattering, transmission loss, and atmospheric illumination effects that cannot be fully captured by appearance transfer alone. Physically grounded atmospheric models therefore remain important for generating realistic environmental degradations and accurately simulating the visibility effects observed in real-world scenes.

\noindent Motivated by these challenges, this thesis investigates atmospheric distribution shift from three complementary perspectives. The first is diffusion-based atmospheric image synthesis and style transfer for synthetic dataset generation. The second explores adaptation strategies for vision-language models under severe atmospheric degradation and source-free deployment constraints. The third investigates physics-based atmospheric synthesis by combining diffusion-based generation, weather-conditioned stylisation, monocular depth estimation, and physically motivated scattering models. Together, these directions provide a holistic framework for the generation, understanding, and mitigation of atmospheric distribution shifts in modern vision systems.

% =====================================================
\section{Problem Statement}
\label{sec:problem_statement}

\noindent Modern vision systems are typically trained on clean, well-curated datasets that do not represent the diversity and severity of environmental conditions encountered in real-world deployments. Atmospheric conditions such as fog, haze, snow, smoke, and sandstorms cause significant variations in image statistics, reductions in visibility, suppression of semantic information, and substantial distribution shifts that deteriorate the performance of vision transformers and vision-language models. Current atmospheric datasets are too small, too homogeneous, and not severe enough to study these effects systematically or to develop robust solutions.

\noindent A further challenge is that many existing style-transfer frameworks are designed to alter global appearance while preserving the semantic content of foreground regions. While such techniques can produce visually compelling atmospheric images, the transformations they generate do not necessarily lead to meaningful changes in the feature representations learned by modern vision models. Perceptual realism alone is therefore insufficient for assessing the quality of synthetic atmospheric generation.

\noindent This thesis tackles three interrelated problems. First, realistic atmospheric imagery generation using controllable diffusion-based frameworks for synthetic dataset creation. Second, the adaptation of modern vision-language models under severe atmospheric distribution shifts, particularly in settings where source-domain data may not be available. Third, the integration of physically grounded atmospheric formation principles into generative pipelines to improve the realism, controllability, and representational impact of synthesised environmental conditions.

\noindent The overall goal is to develop methods for modelling and mitigating atmospheric distribution shifts using generative synthesis, source-free adaptation, and physics-based image formation to improve the robustness of vision systems in real-world environments.

% =====================================================
